% Generated draft: requires mathematical and source review.
\documentclass[11pt]{article}
\usepackage[margin=25mm]{geometry}
\usepackage{fontspec}
\setmainfont{Latin Modern Roman}
\usepackage{amsmath,amssymb,mathtools,unicode-math}
\setmathfont{Latin Modern Math}
\usepackage{enumitem,needspace,xurl,hyperref}
\setlength{\parindent}{0pt}
\setlength{\parskip}{5pt}
\setlength{\emergencystretch}{3em}
\newcommand{\N}{\mathbb{N}}
\newcommand{\Z}{\mathbb{Z}}
\newcommand{\Q}{\mathbb{Q}}
\newcommand{\R}{\mathbb{R}}
\newcommand{\C}{\mathbb{C}}
\newcommand{\F}{\mathbb{F}}
\newcommand{\A}{\mathbb{A}}
\newcommand{\PP}{\mathbb{P}}
\newcommand{\E}{\mathbb{E}}
\newcommand{\eps}{\varepsilon}
\newcommand{\dd}{\,\mathrm{d}}
\begin{document}
\section*{UnsolvedMath ID: 11000294}
\textbf{Source problem code:} catalogue entry
\subsection{Definitions and mathematical statement}
Retain the original cycles of Morita and their generalizations by Morita and by Conant--Vogtmann in the chain models and ranks intended by the source. The record does not supply their formulas, indexing, degrees, or coefficient complexes.

\emph{Question 5.2.} Are the original Morita cycles nontrivial in homology? Are the generalized cycles also nontrivial? For each intended cycle \(z\), nontriviality means
\[
 \partial z=0,\qquad [z]\neq0,
\]
so \(z\) is not a boundary in the relevant chain complex. In the rational-homology formulation suggested by the source, ask for nonvanishing of all intended classes in the appropriate \(H_*(\operatorname{Out}(F_n);\Q)\), with rank and degree attached to each cycle. The source then asks whether these classes account for all rational homology of \(\operatorname{Out}(F_n)\); the precise rank range and whether ``all'' means positive-degree spanning, inclusion of the degree-zero class, or some specified closure of the named families are not stated.

The source motivates completeness by saying no other classes had been found ``to date.'' This is historical, not a current assertion that no other classes exist. The supplied research summary reports the first two Morita classes nonzero for \(\operatorname{Out}(F_4)\) and \(\operatorname{Out}(F_6)\), mentions low-rank \(\operatorname{Aut}(F_n)\) results, and describes higher and graph-parametrized nonvanishing as unproved. It credits Conant--Vogtmann with a one-step stabilization vanishing result in \(\operatorname{Aut}(F_n)\). Those supplied progress claims lack precise citations and maps here; an \(\operatorname{Aut}\)-stabilization statement is not automatically a statement about \(\operatorname{Out}\), and stable vanishing does not establish unstable vanishing.
\subsection{Short English statement}
Determine nonvanishing of Morita's cycles and their generalizations, then whether they exhaust the intended rational homology of free-group outer automorphism groups. The source's historical completeness motivation and reported low-rank/stabilization progress do not settle higher unstable classes.
\subsection{Sources}
\begin{itemize}
\item Local UnsolvedMath registry, permanent ID 11000294, source code catalogue entry, Question 5.2, PDF page 336. Morita, Conant, and Vogtmann are explicitly credited; no individual proposer field is supplied.
\item Farb, \emph{Problems on Mapping Class Groups and Related Topics}. Supplied URL: \url{https://www.math.uchicago.edu/~farb/papers/mcgbook.pdf}; recorded access 2026-07-29; extraction pdf-font-structure. The volume contains 329 explicitly typeset headings across chapters, with repeated labels disambiguated by page.
\item The summary reports first-two-class nonvanishing in ranks \(4,6\), unspecified low-rank \(\operatorname{Aut}\) results, graph-cycle questions, and Conant--Vogtmann stabilization vanishing, without full bibliographic entries or precise statements.
\item Supplied status open and OPEN-TRIAGE. The summary says open in general, while describing those partial results. Current primary-literature status and redistribution terms remain NEEDS\_REVIEW.
\end{itemize}
\end{document}
